\documentclass[11pt,a4paper]{article}
\usepackage[T1]{fontenc}
\usepackage[utf8]{inputenc}
\usepackage{amsmath,amssymb,amsthm}
\usepackage[margin=1in]{geometry}
\usepackage[colorlinks=true,linkcolor=blue,urlcolor=blue]{hyperref}
\theoremstyle{plain}
\newtheorem{theorem}{Theorem}
\newtheorem{lemma}[theorem]{Lemma}
\newtheorem{proposition}[theorem]{Proposition}
\newtheorem{corollary}[theorem]{Corollary}
\theoremstyle{definition}
\newtheorem{remark}[theorem]{Remark}

\title{The empirical recurrence for OEIS A187373, proved by transfer matrix}
\author{Adrian Perez Fontelles\\ \small Independent researcher}
\date{6 September 2026}
\begin{document}
\maketitle

\begin{abstract}
OEIS A187373 counts arrays of width $3$ over $\{0,1,2,3\}$ in which every $3\times3$ subblock
satisfies a condition on its nine entries. The entry carries an empirical recurrence of order
$40$ contributed by R.~H.~Hardin. It is true, and it is decidable rather than empirical: a
$3\times3$ block lies in three consecutive rows, so the state is the last two rows and the
count is a walk count on $S=4161$ states.
\end{abstract}

\noindent\small 2020 Mathematics Subject Classification. 05A15, 05B20, 15A18.\normalsize

\section{The conjecture}

OEIS A187373 is ``Half the number of (n+2)X3 0..3 arrays with each 3X3 subblock having sum 13 or 14.''. It has offset $1$ and begins
\[
30276, 578832, 11193852, 222367896, 4615930512, 96136061832, 2017081911744, 42927291866736,\ \dots
\]
The entry states, as a conjecture contributed by its author and never marked settled:
\begin{quote}
Empirical: a(n)=12*a(n-1)+144*a(n-2)+8364*a(n-3)-85968*a(n-4)-1025856*a(n-5)-22540932*a(n-6)+175623984*a(n-7)+2326114368*a(n-8)+23022501696*a(n-9)-124292765952*a(n-10)-2240639656704*a(n-11)-6241406088960*a(n-12)+3951551070720*a(n-13)+929936095649280*a(n-14)-2199121049454336*a(n-15)+27388235408044032*a(n-16)-126009950324438016*a(n-17)+1363797544592636928*a(n-18)-10200775372373569536*a(n-19)-9861110258169065472*a(n-20)-200171738877536047104*a(n-21)+1251463466276954062848*a(n-22)+2567465691549838442496*a(n-23)+4989615573505839759360*a(n-24)-39261875389490334597120*a(n-25)-112864472074442185113600*a(n-26)+850401008516066652979200*a(n-27)-1287421175527968905625600*a(n-28)-3120464501194565615616000*a(n-29)-63180604151805785407488000*a(n-30)+42574222265676925501440000*a(n-31)+166130233406853532876800000*a(n-32)+1270747183358845491609600000*a(n-33)+425660202835763842252800000*a(n-34)+2189044430502025494528000000*a(n-35)-2240376895000509677568000000*a(n-36)+16131556022908260188160000000*a(n-37)-65768732975904042516480000000*a(n-38)+34116345635915641651200000000*a(n-39)+227442304239437611008000000000*a(n-40)
\end{quote}
As of the ``Last modified'' line on the live entry (Oct 03 2025 03:55:22, revision 10) this is still
recorded as empirical, and nothing on the entry records it as proved.

\section{What is being counted}

The arrays have $3$ columns and a growing number of rows, with entries in $\{0,1,2,3\}$. The
condition is imposed on every $3\times3$ block of the array --- every choice of three
consecutive rows and three consecutive columns --- and asks that the nine entries sum to one of $13,\,14$. The entry publishes $1/2$ of that count; the division is by a constant and changes nothing about the recurrence.



\section{Three rows}

A $3\times3$ block occupies three consecutive rows, so everything the condition asks is decided
by three consecutive rows and nothing wider. Take the array one row at a time and let the state
be the last two rows written, with a distinguished start standing for the array before its
first row. Appending a row completes the blocks whose bottom row is the new one, and the step
is allowed exactly when each of those blocks passes.

Every block of the finished array is tested exactly once, at the step that writes its bottom
row, and no step imposes anything the definition does not ask. The first two rows complete no
block at all, which is why they are laid down without a test. Every state may end an array, so
$\tau=\mathbf{1}$. Thus
\[
a(n)\;=\;\iota^{\!\top}M^{\,j}\tau ,\qquad S=4161,
\]
for the transfer matrix $M$ on those $S$ states. The walk index $j$ and the entry's own $n$
differ by a constant, read off by matching the published terms rather than assumed. In
particular $a$ is $C$-finite.

\section{The proof}

\begin{lemma}\label{lem:crit}
Let $q(t)=t^{40}-\sum_i c_it^{40-i}$ be the characteristic polynomial of the
conjectured recurrence. The residual $a(n)-\sum_ic_ia(n-i)$ equals
$\iota^{\!\top}M^{\,j}q(M)\tau$ for the corresponding walk index $j$, so the recurrence
holds from some point on if and only if $\iota^{\!\top}M^{j}q(M)\tau=0$ for all large $j$,
and it is enough to examine $j<S$.
\end{lemma}

\begin{proof}
Factor $M^{j}$ out of $M^{j+40}-\sum_ic_iM^{j+40-i}$. For the bound, $M^{S}$
is an integer combination of $I,M,\dots,M^{S-1}$ by Cayley--Hamilton, so the numbers
$u_j=\iota^{\!\top}M^{j}q(M)\tau$ obey the monic recurrence given by the characteristic
polynomial of $M$; $S$ consecutive zeros force every later one, and the last nonzero $u_j$
pins the threshold exactly.
\end{proof}

\begin{theorem}
\[
a(n)\;=\;12\,a(n-1) + 144\,a(n-2) + 8364\,a(n-3) - 85968\,a(n-4) - 1025856\,a(n-5) - 22540932\,a(n-6) + 175623984\,a(n-7) + 2326114368\,a(n-8) + 23022501696\,a(n-9) - 124292765952\,a(n-10) - 2240639656704\,a(n-11) - 6241406088960\,a(n-12) + 3951551070720\,a(n-13) + 929936095649280\,a(n-14) - 2199121049454336\,a(n-15) + 27388235408044032\,a(n-16) - 126009950324438016\,a(n-17) + 1363797544592636928\,a(n-18) - 10200775372373569536\,a(n-19) - 9861110258169065472\,a(n-20) - 200171738877536047104\,a(n-21) + 1251463466276954062848\,a(n-22) + 2567465691549838442496\,a(n-23) + 4989615573505839759360\,a(n-24) - 39261875389490334597120\,a(n-25) - 112864472074442185113600\,a(n-26) + 850401008516066652979200\,a(n-27) - 1287421175527968905625600\,a(n-28) - 3120464501194565615616000\,a(n-29) - 63180604151805785407488000\,a(n-30) + 42574222265676925501440000\,a(n-31) + 166130233406853532876800000\,a(n-32) + 1270747183358845491609600000\,a(n-33) + 425660202835763842252800000\,a(n-34) + 2189044430502025494528000000\,a(n-35) - 2240376895000509677568000000\,a(n-36) + 16131556022908260188160000000\,a(n-37) - 65768732975904042516480000000\,a(n-38) + 34116345635915641651200000000\,a(n-39) + 227442304239437611008000000000\,a(n-40)
\]
for every $n>40$.
\end{theorem}

\begin{proof}
The vector $w=q(M)\tau$ was formed by $40$ matrix--vector products in exact integer
arithmetic and $\iota^{\!\top}M^{j}w$ evaluated for $j=0,1,\dots$, again exactly; those
integers vanish from the index corresponding to $n=40+1$ onwards and the run continued
until the working vector was identically zero, which settles every larger $j$ at once.
\end{proof}

\section{Verification}

Four checks, each independent of the algebra above.

First, the model reproduces all $14$ terms the entry publishes, exactly, in integer
arithmetic.

Second, the count was recomputed for the smallest arrays straight from the definition: every
array was written out cell by cell and each $3\times3$ block tested directly on its nine
entries. No rows, no window, no state and no transfer matrix are involved in that computation,
and the published terms came back.

Third, the conjectured recurrence was evaluated directly on the published terms, with no
matrices involved, and holds wherever the proved range and the data overlap.

Fourth, the annihilation test of Lemma~\ref{lem:crit} was carried out in exact integer
arithmetic on the whole vector rather than on a sampled prefix.

\begin{thebibliography}{9}
\bibitem{oeis} The OEIS Foundation, \emph{The On-Line Encyclopedia of Integer
Sequences}, \url{https://oeis.org}, 2026. Sequence A187373.
\bibitem{stanley} R.~P.~Stanley, \emph{Enumerative Combinatorics}, Volume 1, 2nd ed.,
Cambridge University Press, 2011. (Transfer-matrix method, Section 4.7.)
\bibitem{horn} R.~A.~Horn and C.~R.~Johnson, \emph{Matrix Analysis}, 2nd ed., Cambridge
University Press, 2013.
\end{thebibliography}

\end{document}
